\begin{lemma}[The resolvent of a self-adjoint operator, and its comparison with the operator]
Let $E$ be a finite-dimensional real inner-product space, let $B:E\to E$ be a
self-adjoint continuous linear operator, i.e.\
$\langle B\xi,\zeta\rangle=\langle\xi,B\zeta\rangle$ for all $\xi,\zeta\in E$,
and let $\varepsilon>0$ satisfy $\varepsilon\|B\|\le\tfrac12$. Then there is a
continuous linear operator $A_\varepsilon:E\to E$ such that
\begin{enumerate}
\item[(R1)] $A_\varepsilon$ is self-adjoint:
  $\langle A_\varepsilon\xi,\zeta\rangle
   =\langle\xi,A_\varepsilon\zeta\rangle$ for all $\xi,\zeta\in E$;
\item[(R2)] $A_\varepsilon\xi=B\bigl(\xi+\varepsilon A_\varepsilon\xi\bigr)$ for
  every $\xi\in E$;
\item[(R3)] $\bigl(\xi+\varepsilon A_\varepsilon\xi\bigr)
  -\varepsilon B\bigl(\xi+\varepsilon A_\varepsilon\xi\bigr)=\xi$ for every
  $\xi\in E$;
\item[(R4)] $\varepsilon^{-1}\langle\xi,\xi\rangle
  +\langle A_\varepsilon\xi,\xi\rangle>0$ for every $\xi\ne0$;
\item[(R5)] $\|A_\varepsilon\|\le2\|B\|$;
\item[(R6)] for every $\eta\ge2\varepsilon$ and every $\zeta\in E$,
  \[
    \langle A_\varepsilon\zeta,\zeta\rangle
     \le\langle B\zeta,\zeta\rangle+\eta\,|B\zeta|^2 .
  \]
\end{enumerate}
\end{lemma}
\begin{proof}
\emph{Step 1: $\mathrm{id}-\varepsilon B$ is invertible.}
From $\varepsilon>0$ and $\varepsilon\|B\|\le\tfrac12$ we get
$\|\varepsilon B\|=\varepsilon\|B\|\le\tfrac12<1$, so the Neumann series
$T=\sum_{k\ge0}(\varepsilon B)^k$ converges in the operator norm and satisfies
$T(\mathrm{id}-\varepsilon B)=(\mathrm{id}-\varepsilon B)T=\mathrm{id}$; thus
$\mathrm{id}-\varepsilon B$ is invertible with
$(\mathrm{id}-\varepsilon B)^{-1}=T$. Each partial sum of the series is a
polynomial in $B$, hence self-adjoint (a sum of powers of a self-adjoint
operator is self-adjoint) and commuting with $B$; since
$\xi\mapsto\langle S\xi,\zeta\rangle-\langle\xi,S\zeta\rangle$ and
$S\mapsto SB-BS$ are continuous in $S$, the limit $T$ is again self-adjoint and
commutes with $B$. Also
$\|T\|\le\sum_{k\ge0}\|\varepsilon B\|^k\le\sum_{k\ge0}2^{-k}=2$, and more
precisely $\|T\|\le(1-\varepsilon\|B\|)^{-1}\le2$.

Define $A_\varepsilon=BT=B(\mathrm{id}-\varepsilon B)^{-1}$, a continuous linear
operator on $E$.

\emph{(R1).} $B$ and $T$ are self-adjoint and commute, so for all $\xi,\zeta$,
\[
  \langle BT\xi,\zeta\rangle=\langle T\xi,B\zeta\rangle
   =\langle\xi,TB\zeta\rangle=\langle\xi,BT\zeta\rangle .
\]

\emph{The resolvent identity.}
Using $(\mathrm{id}-\varepsilon B)T=\mathrm{id}$ and the distributivity of
composition over addition,
\begin{equation}\label{eq:srr-resolvent}
  \mathrm{id}+\varepsilon A_\varepsilon
   =\mathrm{id}+\varepsilon BT
   =(\mathrm{id}-\varepsilon B)T+\varepsilon BT
   =\bigl[(\mathrm{id}-\varepsilon B)+\varepsilon B\bigr]T
   =T,
\end{equation}
where the second equality replaced $\mathrm{id}$ by
$(\mathrm{id}-\varepsilon B)T$.

\emph{(R2).} By \eqref{eq:srr-resolvent}, $T\xi=\xi+\varepsilon A_\varepsilon\xi$
for every $\xi$, so
$A_\varepsilon\xi=B(T\xi)=B\bigl(\xi+\varepsilon A_\varepsilon\xi\bigr)$.

\emph{(R3).} Again by \eqref{eq:srr-resolvent} and
$(\mathrm{id}-\varepsilon B)T=\mathrm{id}$,
\[
  \bigl(\xi+\varepsilon A_\varepsilon\xi\bigr)
   -\varepsilon B\bigl(\xi+\varepsilon A_\varepsilon\xi\bigr)
   =(\mathrm{id}-\varepsilon B)(T\xi)=\xi .
\]

\emph{(R4).} For $\xi\ne0$, the Cauchy--Schwarz inequality and the definition of
the operator norm give $\langle B\xi,\xi\rangle\le|B\xi|\,|\xi|\le\|B\||\xi|^2$,
hence
\[
  \langle(\mathrm{id}-\varepsilon B)\xi,\xi\rangle
   =|\xi|^2-\varepsilon\langle B\xi,\xi\rangle
   \ge(1-\varepsilon\|B\|)|\xi|^2\ge\tfrac12|\xi|^2>0 .
\]
So $\mathrm{id}-\varepsilon B$ is self-adjoint and positive definite. Its
inverse $T$ is then positive definite as well: for $x\ne0$ we have $Tx\ne0$
(because $T$ is injective, being invertible), and
\[
  \langle Tx,x\rangle=\bigl\langle Tx,(\mathrm{id}-\varepsilon B)Tx\bigr\rangle
   =\bigl\langle(\mathrm{id}-\varepsilon B)(Tx),Tx\bigr\rangle>0,
\]
using $x=(\mathrm{id}-\varepsilon B)Tx$ and the self-adjointness of
$\mathrm{id}-\varepsilon B$. Consequently, by \eqref{eq:srr-resolvent},
\[
  \varepsilon^{-1}\langle\xi,\xi\rangle+\langle A_\varepsilon\xi,\xi\rangle
   =\varepsilon^{-1}\bigl\langle(\mathrm{id}
     +\varepsilon A_\varepsilon)\xi,\xi\bigr\rangle
   =\varepsilon^{-1}\langle T\xi,\xi\rangle>0\qquad(\xi\ne0).
\]

\emph{(R5).} $\|A_\varepsilon\|=\|BT\|\le\|B\|\,\|T\|\le2\|B\|$ by Step 1.

\emph{(R6).} Since $B$ is self-adjoint on the finite-dimensional space $E$, the
spectral theorem supplies an orthonormal basis $(v_i)_{i=1}^n$ of $E$ with
$Bv_i=b_iv_i$ and $b_i\in\R$; moreover
$|b_i|=|Bv_i|\le\|B\|\,|v_i|=\|B\|$. For each $i$,
$(\mathrm{id}-\varepsilon B)v_i=(1-\varepsilon b_i)v_i$ with
$1-\varepsilon b_i\ge1-\varepsilon\|B\|\ge\tfrac12>0$, so $Tv_i=
(1-\varepsilon b_i)^{-1}v_i$ and therefore $A_\varepsilon v_i=a_iv_i$ with
$a_i=b_i/(1-\varepsilon b_i)$. Now
\[
  a_i-b_i=\frac{b_i-b_i(1-\varepsilon b_i)}{1-\varepsilon b_i}
   =\frac{\varepsilon b_i^2}{1-\varepsilon b_i}\le2\varepsilon b_i^2\le\eta b_i^2 ,
\]
using $0<1-\varepsilon b_i$, $1-\varepsilon b_i\ge\tfrac12$ (so
$(1-\varepsilon b_i)^{-1}\le2$), $\varepsilon b_i^2\ge0$ and
$2\varepsilon\le\eta$. Writing $\zeta=\sum_ic_iv_i$ with $c_i\in\R$ and using
orthonormality,
\[
  \langle A_\varepsilon\zeta,\zeta\rangle=\sum_ia_ic_i^2,\qquad
  \langle B\zeta,\zeta\rangle=\sum_ib_ic_i^2,\qquad
  |B\zeta|^2=\sum_ib_i^2c_i^2 .
\]
Multiplying the termwise inequality $a_i-b_i\le\eta b_i^2$ by $c_i^2\ge0$ and
summing over $i$ gives
$\langle A_\varepsilon\zeta,\zeta\rangle-\langle B\zeta,\zeta\rangle
\le\eta|B\zeta|^2$, which is (R6).
\end{proof}
